% © 2026 Ing. David Jaroš — CC BY-NC-ND 4.0
\documentclass[11pt,a4paper]{article}
\usepackage[margin=2.5cm]{geometry}
\usepackage{amsmath,amssymb,amsthm,mathtools}
\usepackage[most]{tcolorbox}
\newtheorem{theorem}{Theorem}
\newtheorem{corollary}[theorem]{Corollary}
\newtheorem{remark}[theorem]{Remark}
\newcommand{\dd}{\ddagger}
\title{GAP-10T Minimal One-Connection GR Audit:\\
An Architecture-Level No-Go and the Exact Completion Fork}
\author{Ing. David Jaro\v{s}}
\date{26 July 2026}
\begin{document}\maketitle

\begin{abstract}
We combine two previously proved UBT results: the concurrent-vector obstruction
for a torsion-free tetrad generated by the pure Lorentz pair, and the local
right inverse for an arbitrary tetrad obtained by adding composite contortion.
The combination yields an architecture-level dichotomy. If the one connection
inside $D\Theta$ is also the physical metric connection and the physical branch
is exact torsion-free general relativity, then the generated tetrad admits a
concurrent vector and generic GR geometries, including the non-flat
Schwarzschild exterior, are excluded. If nonzero contortion is retained to
represent arbitrary tetrads, the same physical connection has nonzero torsion
and the theory is not exact torsion-free GR. Therefore the minimal architecture
with one and the same physical Lorentz connection cannot simultaneously provide
universal local single-$\Theta$ representability and exact generic GR. The
result does not exclude UBT gravity; it proves that an exact-GR completion must
relax at least one assumption, most economically by separating a nonpropagating
jet connection from the physical Levi--Civita connection, or else by accepting
a genuinely torsionful gravity theory.
\end{abstract}

\section{Assumptions}
Let
\begin{equation}
 E_\mu=\mathcal N_0^{-1/2}D^{\Omega}_\mu\Theta,
 \qquad
 D^{\Omega}_\mu X=\partial_\mu X+\Omega_\mu X+X\Omega_\mu^{\dd},
 \label{eq:generated}
\end{equation}
with $E_\mu$ a nondegenerate Lorentz tetrad.  The minimal
\emph{one-connection} interpretation imposes:
\begin{enumerate}
 \item $\Omega$ is the only Lorentz connection in the defining jet;
 \item the same $\Omega$ is the spin lift of the physical metric-compatible
 affine connection that enters curvature and matter transport;
 \item no independent relative left/right connection acts on $\Theta$.
\end{enumerate}
Exact ordinary GR further requires the physical connection to be the
Levi--Civita connection of $E$ on the spinless vacuum branch.

\section{The dichotomy}
\begin{theorem}[Minimal one-connection exact-GR no-go]
\label{thm:main}
Under the assumptions above, the same one-connection architecture cannot have
both of the following properties:
\begin{enumerate}
 \item every smooth Lorentzian tetrad is locally representable by a single
 $\Theta$ through Eq.~\eqref{eq:generated};
 \item the physical vacuum branch is exact torsion-free GR for generic metrics.
\end{enumerate}
\end{theorem}

\begin{proof}
If the physical branch is torsion-free, metric compatibility fixes
$\Omega=\mathring\Omega(E)$.  The previously proved concurrent-vector theorem
then applies to the Lorentz-real projection of $\Theta$: there is a spacetime
vector $V$ such that
\begin{equation}
 \mathring\nabla_\mu V^\nu=\delta_\mu{}^\nu,
 \qquad
 \mathcal L_V g_{\mu\nu}=2g_{\mu\nu},
 \qquad
 R^\rho{}_{\sigma\mu\nu}V^\sigma=0.
 \label{eq:concurrent}
\end{equation}
Generic torsion-free Einstein geometries do not admit this structure; in
particular, the non-flat Schwarzschild vacuum exterior with $M\ne0$ admits no
proper homothety of Eq.~\eqref{eq:concurrent}.  Hence exact generic GR fails.

Conversely, the local representer theorem constructs a right inverse for every
smooth tetrad by using
\begin{equation}
 \Omega=\mathring\Omega(E)+K(E,\Theta),
 \qquad K\ne0\ \hbox{in general}.
\end{equation}
For a metric-compatible connection, contortion and torsion are in one-to-one
algebraic correspondence.  Therefore $K\ne0$ makes the same physical affine
connection torsionful.  It may define an Einstein--Cartan or other modified
branch, but it is not the Levi--Civita connection of exact ordinary GR.
Thus the two desired properties cannot hold simultaneously under the same
one-connection interpretation.
\end{proof}

\begin{corollary}[What must be relaxed]
An exact generic-GR completion of the single-$\Theta$ tetrad map must relax at
least one of the following:
\begin{enumerate}
 \item identify the jet connection with the physical spacetime connection;
 \item restrict the derivative to the pure Lorentz pair of
 Eq.~\eqref{eq:generated};
 \item require exact torsion-free GR rather than a torsionful modification;
 \item require local representability of an arbitrary tetrad by one $\Theta$.
\end{enumerate}
\end{corollary}

\section{Controlled completion fork}
The smallest exact-GR architecture suggested by the theorem is
\begin{equation}
 \widehat\Omega=\mathring\Omega(E)+K_J[E,\Theta],
 \qquad
 E_\mu=\mathcal N_0^{-1/2}D^{\widehat\Omega}_\mu\Theta,
 \label{eq:jet}
\end{equation}
for the defining jet, while physical curvature and ordinary matter transport
use
\begin{equation}
 \Omega_{\rm phys}=\mathring\Omega(E),
 \qquad T_{\rm phys}=0.
 \label{eq:phys}
\end{equation}
A companion theorem gives an explicit local composite Lorentz plus relative-
central jet right inverse for every non-null Lorentz-real projection of
$\Theta$.  It must not be interpreted as physical spacetime torsion.  The
remaining issue is dynamical selection of that right inverse and of the tetrad,
not local kinematic existence.

The alternative is to keep one connection and interpret the result as a
physical torsionful theory.  That route requires observational bounds and an
action that dynamically selects the composite contortion.

\section{What remains for a complete GR derivation}
Even after choosing the split exact-GR architecture, two independent lemmas
remain:
\begin{enumerate}
 \item \textbf{Jet dynamics lemma:} derive from the canonical action the
 selection of the explicit split-jet right inverse and of $E[\Theta]$, and
 prove that the relative central/Lorentz jet components are nonpropagating.
 \item \textbf{Curvature-action lemma:} derive the Hilbert--Palatini/Einstein--
 Hilbert term, including its sign and Newton coefficient, from the canonical
 UBT action.  The locked quadratic kinetic scalar is only a volume term and
 cannot supply $R$ by algebraic rewriting.
\end{enumerate}
Once these are proved, standard Palatini variation and the already proved
connection-reconstruction theorem give the Einstein--$\Lambda$ equations on
the torsion-free physical branch.  Until then, the endpoint is a controlled
conditional completion rather than an unconditional derivation from the
original minimal action.

\begin{tcolorbox}[colback=green!6!white,colframe=green!55!black,
 title=Exact status]
\textbf{GAP-10T-MINIMAL-ONE-CONNECTION-GR: CLOSED AS NO-GO [L1].}
The same one metric-compatible physical Lorentz connection cannot provide both
universal local single-$\Theta$ representability and exact generic torsion-free
GR.

\medskip
\textbf{GAP-10T-JET-KIN: CLOSED LOCALLY [L1].}
The split exact-GR architecture has an explicit local composite jet right
inverse while physical curvature remains Levi--Civita.

\medskip
\textbf{GAP-10T-JET-DYN: OPEN.}
The canonical action has not selected the tetrad or proved nonpropagation of the
jet completion.

\medskip
\textbf{GAP-10D: NARROWED, NOT CLOSED.}
The physical Einstein endpoint is conditionally standard once its curvature
term is present; the canonical origin and coefficient of that term remain to
be derived.
\end{tcolorbox}

\begin{thebibliography}{9}
\bibitem{wald1984} R.~M.~Wald, \emph{General Relativity}, University of Chicago Press (1984).
\bibitem{hehl1976} F.~W.~Hehl, P.~von der Heyde, G.~D.~Kerlick, and J.~M.~Nester,
``General Relativity with Spin and Torsion: Foundations and Prospects,''
\emph{Rev. Mod. Phys.} \textbf{48}, 393--416 (1976).
\end{thebibliography}
\end{document}
